\chapter{Positionnement et discussion}
\label{ch:discussion}

Ce chapitre prend du recul sur le travail réalisé : il récapitule les notions de
plus court chemin mises en œuvre, revient sur un algorithme essayé puis écarté
(IDA*), et met en lumière un parallèle entre les landmarks d'ALT et les Pattern
Databases.

\section{Récapitulatif des notions couvertes}

Il est d'abord utile de dresser le bilan de ce qui a été traité. Le
tableau~\ref{tab:mapping} récapitule les notions fondamentales de plus
court chemin abordées et renvoie, pour chacune, à l'endroit du rapport qui la
met en œuvre.

\begin{table}[htbp]
  \centering
  \small
  \begin{tabular}{@{}ll@{}}
    \toprule
    Notion & Où elle apparaît \\
    \midrule
    Recherche en meilleur d'abord, espace d'états & ch.~\ref{ch:modelisation} \\
    Parcours en largeur (BFS) / profondeur (DFS)   & \S\ref{sec:baselines}, \S\ref{sec:bench} \\
    Greedy Best-First Search                       & \S\ref{sec:greedy} \\
    Dijkstra                                       & \S\ref{sec:dijkstra} (référence) \\
    A*, $f = g + h$, heuristique admissible        & \S\ref{sec:astar} \\
    « A* généralise Dijkstra »                     & A* avec $h \equiv 0$ \\
    Heuristique triangulaire / ALT                 & \S\ref{sec:alt} (étude principale) \\
    Optimisations de structures (tableau de piles) & \S\ref{sec:bucket}, \S\ref{sec:bucket-result} \\
    Landmarks actifs                               & \S\ref{sec:active}, \S\ref{sec:bench} \\
    \bottomrule
  \end{tabular}
  \caption{Notions de recherche de plus court chemin et leur traitement dans le
  projet.}
  \label{tab:mapping}
\end{table}

À ces fondamentaux, le projet ajoute : la recherche bidirectionnelle (Dijkstra et
ALT à potentiel moyen), les \textbf{Arc Flags} et les \textbf{Contraction
Hierarchies}, la sélection \texttt{avoid} des landmarks, une \textbf{étude de
passage à l'échelle} prouvant que les gains croissent avec la taille
(\S\ref{sec:scaling}), et une \textbf{analyse statistique} de la difficulté des
requêtes (\S\ref{sec:distribution}).

\section{IDA* : implémenté, testé, puis écarté}
\label{sec:idastar}

Par souci d'honnêteté sur la démarche, signalons que l'algorithme \textbf{IDA*}
(\emph{Iterative Deepening} A*, \citealp{korf1985ida}) a lui aussi été implémenté
et testé pendant le projet, avant d'être écarté. IDA* est conçu pour les espaces d'états
\emph{implicites et exponentiels} sous forte contrainte mémoire (typiquement les
puzzles) : il remplace les listes ouverte et fermée d'A* par une succession de
parcours en profondeur bornés par un seuil croissant sur $f = g + h$, ne gardant
en mémoire que le chemin courant (mémoire linéaire au lieu de linéaire en nombre
de nœuds visités).

Sur un graphe routier \emph{explicite}, qui tient aisément en mémoire, ce
compromis est contre-productif. Faute de liste fermée, IDA* \textbf{ré-explore}
les mêmes nœuds à chaque nouveau seuil ; or, sur un graphe comportant de nombreux
cycles et des poids non uniformes, les valeurs distinctes de $f$ sont
innombrables, si bien que le nombre de ré-expansions explose. IDA* rend alors
exactement la même route qu'A*, mais bien plus lentement et sans le moindre gain.
Nous l'avons donc écarté au profit d'A* et d'ALT, qui règlent chaque nœud une
seule fois, cohérent avec le fait que, ici, la difficulté est l'échelle du
graphe et non la mémoire.

\section{Landmarks $=$ Pattern Databases pour graphe explicite}
\label{sec:pdb-alt}

Nous terminons par un rapprochement qui éclaire la nature même d'ALT. La borne
par landmarks, telle qu'on l'a introduite, pourrait passer pour une astuce
propre au routage ; elle est en réalité l'instance, sur un graphe explicite,
d'un principe général de la recherche heuristique, illustré par les
\textbf{Pattern Databases} (PDB) utilisées pour les puzzles
\citep{culberson1998pdb}. Le signaler n'est pas un simple parallèle décoratif :
cela explique \emph{pourquoi} l'heuristique ALT fonctionne et la relie à une
idée éprouvée, ce qui donne du recul sur le choix de conception.

\medskip

Les deux approches répondent en effet à la même idée : \emph{pré-calculer des
distances exactes vers un sous-ensemble d'ancres, puis en déduire par relaxation
une borne inférieure admissible}. La table~\ref{tab:pdb-alt} les met en regard.

\begin{table}[htbp]
  \centering
  \small
  \begin{tabular}{@{}lll@{}}
    \toprule
    & Pattern Database & Landmarks (ALT) \\
    \midrule
    Espace d'états      & implicite, exponentiel      & explicite (graphe en mémoire) \\
    Objet pré-calculé   & distances exactes vers un état projeté & distances exactes vers $k$ nœuds \\
    Borne admissible    & distance dans l'abstraction & inégalité triangulaire \\
    Combinaison         & max (ou somme, additives)   & max sur les landmarks \\
    Utilisée par        & IDA* / A*                   & A* \\
    \bottomrule
  \end{tabular}
  \caption{Le même principe, deux incarnations selon la nature de l'espace d'états.}
  \label{tab:pdb-alt}
\end{table}

Une PDB projette l'état sur un sous-problème (par exemple un sous-ensemble de
tuiles du taquin), résout exhaustivement ce sous-problème par analyse rétrograde,
et utilise la distance dans cette abstraction comme borne inférieure admissible.
ALT fait l'analogue direct pour un graphe explicite : il projette la distance
« vers $t$ » sur ce qu'en révèlent quelques ancres $L_i$, via
$d(u,t) \ge |d(u,L_i) - d(t,L_i)|$. Dans les deux cas, on paie un pré-calcul
unique pour obtenir une heuristique bien plus forte que l'estimation géométrique
de base, et dans les deux cas on combine plusieurs bornes par un maximum. Cette
correspondance montre que le projet, bien qu'il ne code pas de PDB, en incarne
pleinement le principe, adapté à la structure de son problème.
